\SourceLeaf{333}

\[
\sin(-1)\sin(2+3)\sin(2+5) = \sin(-2)\sin(1+4)\sin(1+6),
\]

wobei die Größen

\[
-1, \qquad 2+3, \qquad 2+5; \qquad -2, \qquad 1+4, \qquad 1+6
\]

denselben Bedingungen genügen, die für die Größen

\[
+1, \qquad 3, \qquad 5; \qquad +2, \qquad 4, \qquad 6
\]

durch die Gleichungen (1), (2) ausgedrückt waren. Daher lassen sich auch alle Folgerungen übertragen.

\begin{center}
[6.]

[\emph{Handbuch}, S.\ 278 bis 281.]

\textbf{NOCH ZUR BEURTEILUNG DER PHYSISCHEN MÖGLICHKEIT DER DATEN IN POTHENOTS AUFGABE.}
\end{center}

$a$, $b$, $c$, $A$, $B$, $C$, $\alpha$, $\beta$, $\gamma$, $\mathfrak{A}$, $\mathfrak{B}$, $\mathfrak{C}$ haben dieselbe Bedeutung wie oben. Im Einzelnen

\begin{center}\begin{tabular}{@{}lp{9.8cm}@{}}
$a,b,c$ & Seiten;\\
$A,B,C$ & Winkel des Dreiecks aus den gegebenen Punkten (1), (2), (3), mit dem Radius $=r$;\\
$\alpha,\beta,\gamma,\rho$ & Entfernungen des gesuchten Punkts (0) von $A,B,C$ und vom Mittelpunkt dieses Kreises;\\
$\mathfrak{A}$ & Winkel von $02$ nach $03$, ebenso gemessen wie $A$ von $12$ nach $13$ (d.\,h. beide von rechts nach links oder beide umgekehrt), ebenso $\mathfrak{B},\mathfrak{C}$.
\end{tabular}\end{center}

Dann gilt, wie bereits auf p.\ 270 [des Handbuchs] bemerkt wurde,

\[
\frac{\sin(\mathfrak{A}-A)}{a\alpha} = \frac{\sin(\mathfrak{B}-B)}{b\beta} = \frac{\sin(\mathfrak{C}-C)}{c\gamma} = \frac{rr-\rho\rho}{2\alpha\beta\gamma r},
\]

woraus folgt, dass je nachdem, ob (0) innerhalb oder außerhalb dieses Kreises liegt,

\[
\sin(\mathfrak{A}-A), \qquad \sin(\mathfrak{B}-B), \qquad \sin(\mathfrak{C}-C)
\]

alle positiv oder negativ werden.
